% Your comments for the king and pawn positions. Format: notes/README.md.

% Diagram 1 of the 1500-1599 book: Lakusic - Radosevic, 65th ch-MNE 2013
@ 8/8/8/2k5/8/1K6/2P5/8 w - - 4 78 | c4?
The premature pawn thrust lets Black defend the \vocab{key-squares}{key squares}{} in front of the pawn: after 78...\K c6 it is only a draw.\index{King and pawn endings!Key squares}
Instead, 78.\K c3! takes the \vocab[king and pawn endings]{opposition}{opposition}{The kings stand on the same file (or rank) with one square between them, and the other side is to move, so its king has to give way.}.\index{King and pawn endings!Opposition} Black must give way, and the white king goes round the black one, \vocab[king and pawn endings]{outflanking}{outflanking}{Going round the side of the enemy king when it has given way, to reach the key squares.} it: 78...\K d5 79.\K b4 or 78...\K d6 79.\K d4. White wins.\index{King and pawn endings!Outflanking}

% Diagram 2: Olsen - Gabrielsen, TV2 Elo Open 2017
@ 8/1p6/8/2k5/8/2K5/8/8 b - - 5 64 | b5?
Again a premature pawn thrust. 64...\K b5 or 64...b6 would have won.

@ 8/8/8/1pk5/8/2K5/8/8 w - - 0 65 | Kc2?
King and pawn endings are usually decided by \vocab[king and pawn endings]{zugzwang}{zugzwang}{The side to move would rather pass: every move makes its position worse.}.\index{King and pawn endings!Zugzwang}
65.\K c2? lets Black choose between 65...\K b4 and the \vocab{opposition}{opposition}{} with 65...\K c4. Which of the two is right does not matter to White. The calculation is not difficult, so Black will probably find the right move, and as long as one of the two wins, White loses.
65.\K b3! was correct: it leaves Black no such choice, and the game is a draw.\par
This is worth remembering beyond pawn endings: give your opponent no \vocab[practical play]{easy-choice}{easy choice}{A choice between a few moves, one of which wins, that the opponent can settle without much calculation. Avoid offering one. You can find the move that leaves no such choice without calculating which option would win, and so save your energy for the decisions that really need calculation.}.\index{Practical play!No easy choice} To find 65.\K b3 you do not need to work out whether 65...\K b4 or 65...\K c4 wins. It is enough to see that 65.\K c2 lets Black choose. A decision made without calculation saves energy for the moments that need it.

@ 8/8/8/8/8/1pk5/8/1K6 b - - 1 70 | b2
\diagram{8/8/8/8/8/2k5/1p6/1K6 w - - 0 71}
White is in zugzwang: the king has to leave b1, and the pawn queens. With Black to move, Black would be in zugzwang too: nothing wins, and 71...\K b3 is stalemate. Such a position, where whoever has to move loses something, is a \vocab[king and pawn endings]{mutual-zugzwang}{mutual zugzwang}{A position in which the side to move, whichever it is, would rather pass.}.\index{King and pawn endings!Zugzwang!mutual}

@ 8/8/8/5pk1/8/6K1/8/8 w - - 2 47 | Kg2?
No comment.

@ 8/8/5K2/k6p/8/8/7P/8 b - - 0 63 | Kb6?
White can play for a win: Black's pawn will soon fall. To win, White has to take it and then push
his own pawn.

Black's defence is to use his pawn to lure the white king away, and so gain time for his own king.
Only 63...h4! draws, for example 63...h4 64.\K g5 h3 65.\K h4 \K b5 66.\K xh3 \K c6 67.\K g4 \K d7
68.\K f5 \K e7 69.\K g6 \K f8.

Once the black king reaches f8, White can make no progress. With a rook's pawn the defending king
only has to get to the corner in front of it, h8.\index{King and pawn endings!Rook's pawn} If the
white king gets there first to keep it out, the black king shuttles between f7 and f8. White can
push his pawn, but in the end he is stalemated on h8:\index{King and pawn endings!Stalemate}
\diagram{7K/5k1P/8/8/8/8/8/8 b - - 0 75}
75...\K f8, and White has no legal move. Note that this is the only move: any other lets the white
king out, and the pawn queens.
